\vpage{1}{One root is not the whole problem}{READER'S GUIDE}{Main text: §§1--3, 8--11. Current status: the later scalar/Mellin continuation.}
Paper X separates three questions that can look deceptively similar. Does a residual object exist? Does it extend coherently through every precision? Does the resulting object occur in the required carrier? A successful answer to one question does not supply the next.
\vfig{01_three_questions}
\head{Three models to keep in view}
The scalar equation $x^2+1=0$ has the root $1$ modulo two, but no root modulo four. Over powers of five, the residues $2,7,57,182$ form the beginning of a compatible branch. Neither fact constructs an automorphic carrier. In a supplied finite carrier, an idempotent such as $\diag(1,0)$ detects an occupied subspace; it does not explain where that carrier came from.

Read the appendix in three passes. Pages \pagerefV{3}--\pagerefV{11} explain residual digits, finite counts, and obstruction spaces. Pages \pagerefV{12}--\pagerefV{20} explain directed transport, common data, and finite occupancy. Pages \pagerefV{21}--\pagerefV{29} give small exact models for the newer arithmetic continuation and its remaining limits.
\begin{insight}
\textbf{Current formal boundary.} The latest recovered continuation still records \textbf{17 checked claims and three partial claims: 9.1, 9.2, and 11.1}. This appendix changes none of those statuses. In particular, prescribed Frobenius occurrence is not proved by the finite character calculations below.
\end{insight}
\captiontext{The figures are exact-model illustrations, not pictures of an entire arithmetic moduli space. New Python checks support these examples. Saved Lean evidence is preserved separately; no new kernel run is claimed.}

\vpage{2}{An atlas of the objects}{NOTATION BEFORE PICTURES}{Main text: Definitions 1.1, 5.1, 6.1, 7.1--7.2, and 8.1.}
Keep the coefficient ring, its reduction, and the carrier separate. A residue class is a finite-precision address. A vector in a carrier is a different kind of object, even when both are selected by something called a packet.
\par\noindent\begin{tabularx}{\linewidth}{@{}p{.19\linewidth}YY@{}}
\toprule
Object & What it records & What it does not supply\\\midrule
$\OO,\varpi,k$ & Complete valuation ring, uniformizer, and residue field $k$ of size $q$. & A represented deformation functor.\\
$\bar x$ & One residual object in the declared moduli problem. & Other residual roots or their transports.\\
$X_m(\bar x)$ & Admissible objects modulo $\varpi^m$ reducing to $\bar x$. & A lift of every individual member.\\
$(x_m)_m$ & A branch satisfying $x_{m+1}\bmod\varpi^m=x_m$. & A global object without the stated reconstruction hypotheses.\\
$w_i$ & A coordinate constrained to begin in $\varpi^{w_i}\OO$. & A classification of all singular deformation rings.\\
$Z(T)$ & The formal series $\sum_{m\ge1}\#X_mT^m$. & The cardinality of the inverse limit.\\
$C^\bullet,d$ & Cochains and a differential with $d^2=0$. & Its arithmetic identification with a Galois problem.\\
$G$ & The graph of supplied admissible residual transports. & A proof that a proposed arithmetic edge exists.\\
$\mathscr C$ & Common algebraic polynomials and declared local data. & A member realizing those data at every prime.\\
$C\bm m=\bm b$ & A finite trace system for carrier multiplicities. & Integrality, positivity, or available capacity automatically.\\
$e_TV$ & The joint eigenspace inside an existing finite carrier $V$. & Nonzero occupancy for every desired joint label.\\\bottomrule
\end{tabularx}
\head{Three indices that must not be conflated}
The depth $m$ counts retained digits. A coordinate index $i$ selects a parameter or a packet. A coefficient place $\lambda$ selects a different local field. Increasing $m$ at one $\lambda$ neither changes $i$ nor constructs the member at another $\lambda$.
\takeaway{Specify the ambient object before applying a projector, a reduction, a rank, or a counting formula. The same-looking symbol can refer to quite different spaces.}

\vpage{3}{Higher digits cannot change the root}{RESIDUAL PARTITION}{Main text: Theorems 2.1--2.2 and Corollary 2.3.}
Work in $\ZZ/27\ZZ$. Every element has a unique form $x=3q+r$ with $r\in\{0,1,2\}$ and $q\in\{0,\ldots,8\}$. The row below is the residual root. The column contains all the remaining information at this precision.
\vfig{03_residual_grid}
\head{Read the operation, not just the arrow}
An operation $x\mapsto x+3h$ preserves $r$, because its correction vanishes modulo three. Repeated corrections of this form can visit every column in a row, but can never enter another row. By contrast, adding one changes the residual class. It is therefore not a positive-depth perturbation in this model.

For example, the addresses $1,4,7,10,\ldots,25$ may be deeply connected by higher-digit changes. They remain disjoint from the addresses reducing to zero. In a genuine deformation problem, the same conclusion applies to operations whose reduction is the identity, not to every operation that happens to be called a deformation.
\[
\operatorname{red}(x+3h)=\operatorname{red}(x),
\qquad
\mathcal M(A)=\bigsqcup_{\bar x}\mathcal M_{\bar x}(A).
\]
\takeaway{Vertical completeness can fill one entire residual fiber without reaching a single point in a different fiber.}
\captiontext{All 27 addresses and their finite corrections are checked by integer reduction. Passing to an infinite perturbation stream additionally requires the convergence condition recorded in the formal no-crossing statement.}

\vpage{4}{A branch fixes digits permanently}{INVERSE-LIMIT RECONSTRUCTION}{Main text: Theorem 3.2 and the compatible-digit example in §3.}
The rational number $1/4$ is a five-adic integer. Its normalized residues are obtained by solving $4x_m\equiv1\pmod{5^m}$, not by taking decimal approximations. At the first six depths they are $4,19,94,469,2344,11719$.
\vfig{04_digits}
\head{Compatibility is stronger than independent solvability}
The residues satisfy both
\[
x_{m+1}\equiv x_m\pmod{5^m},
\qquad 4x_m-1=3\cdot5^m.
\]
The latter identity shows that the residual error tends to zero in the five-adic norm. The stabilized digit sequence is
\[
\frac14=4+3\cdot5+3\cdot5^2+3\cdot5^3+\cdots\quad\text{in }\ZZ_5.
\]
This series is not convergent in the ordinary real norm. The topology is part of the statement. To view the branch as a positive-depth coordinate above the residual root $4$, use $t=x-4\in5\ZZ_5$.

An inverse-limit point is the entire compatible family, not its last displayed approximation. The six rows are a finite window; the displayed exact formula supplies a branch at every depth in this particular example. In the represented deformation theorem, compatibility reconstructs maps on all ring elements, with multiplication, continuity, and locality retained.
\takeaway{A later approximation may append a digit. It may not revise a digit already certified by reduction.}

\vpage{5}{A smooth chart has independent choices}{ACTUAL FINITE POINTS}{Main text: Proposition 3.4 and Theorem 6.2.}
Take $R=\ZZ_3[[t_1,t_2]]$ and coefficient-preserving local maps to $\ZZ/27\ZZ$. Locality forces both parameters into the maximal ideal. Each therefore has nine possible values, $0,3,6,\ldots,24$.
\vfig{05_smooth_grid}
\head{Why the exponent is $d(m-1)$}
At depth $m$, the image of the maximal ideal contains $3^{m-1}$ elements. Two independent parameters give $3^{2(m-1)}$ points. The grid displays all $81$ points at depth three, not a random sample. Reduction to depth two maps each coordinate to one of $0,3,6$, and every parent has nine children.
\[
\#X_m=3^{2(m-1)},\qquad
Z(T)=\sum_{m\ge1}\#X_mT^m=\frac{T}{1-9T}.
\]
Counting all $27^2$ parameter assignments at depth three would ignore locality. Counting only tangent directions would describe the first layer but not all later layers. The power-series presentation explains why these choices stay independent through every precision.

This is a coordinate model with a proved presentation. A general deformation ring can impose equations among the parameters, changing the counts and causing branches to die. The next pages retain the same depth convention while changing the allowed coordinates or adding a singular equation.
\takeaway{The smooth count is not a universal rule for deformation spaces. Its independence comes from the actual formal power-series chart.}

\vpage{6}{Weights delay visibility}{WEIGHTED DYADIC CHARTS}{Main text: Proposition 5.3, Theorem 6.3, and §12.2.}
Over $\ZZ_2$, require three coordinates to lie in $2\ZZ_2$, $2^2\ZZ_2$, and $2^4\ZZ_2$. At depth $m$, coordinate $i$ has $\max(m-w_i,0)$ free binary digits. A zero in the table means only one allowed residue: zero.
\vfig{06_weights}
\head{The two different thresholds}
A coordinate of weight $w_i$ first becomes nontrivial when $m=w_i+1$. By contrast, the first separation depth of two distinct fixed values $a,b$ is $v_2(a-b)+1$. Thus $1$ and $-1$ separate at depth two, because their difference has valuation one. Neither threshold is simply the printed coordinate label.

For weights $(1,2,4)$,
\[
N(m)=2^{\max(m-1,0)+\max(m-2,0)+\max(m-4,0)}.
\]
The first counts are $1,2,8,32,256,2048$. Once all three coordinates are active, every additional precision digit multiplies the count by eight. Separating the initial levels from that geometric tail gives
\[
Z(T)=\frac{T(1-6T-8T^2-32T^3)}{1-8T}.
\]
Negative coefficients in the numerator do not mean negative numbers of points; expanding the whole rational expression gives positive coefficients.
\takeaway{Weights change the finite-depth delay. After activation, this chart still has three independent digit streams.}

\vpage{7}{Many approximate roots, one limit}{SINGULAR POWER SEEDS}{Main text: Theorem 6.4 and Corollary 6.6.}
Now impose $x^2=0$ modulo $3^m$. The condition is exactly $2v_3(x)\ge m$, so the roots are the multiples of $3^{\lceil m/2\rceil}$ in that quotient. Their number is $3^{m-\lceil m/2\rceil}$.
\vfig{07_singular_forest}
\head{Count vertices and count branches separately}
There are three roots modulo nine: $0,3,6$. But the only roots modulo $27$ are $0,9,18$, and all three reduce to zero modulo nine. The other two vertices at the previous level have no children. Counting them as permanent branches would be incorrect.
\[
\#X_m=1,3,3,9,9,27,\ldots,
\qquad Z(T)=\frac{T+3T^2}{1-3T^2}.
\]
The count-growth exponent is $\lim_m\log_3(\#X_m)/m=1/2$. Nevertheless, a compatible branch reconstructs $x\in\ZZ_3$ with $x^2=0$. Since $\ZZ_3$ has no nonzero nilpotent, the only such branch is the zero branch.

More generally, $x^e=0$ has count $q^{m-\lceil m/e\rceil}$ and count-growth exponent $1-1/e$. These exponents describe finite thickenings. They do not assert that the set of complete roots has a fractional geometric dimension or more than one point.
\takeaway{The seed zeta function counts finite-level solutions. It is not a generating function for distinct infinite branches.}

\vpage{8}{One equation, several local outcomes}{A SCALAR COMPANION MODEL}{Main text: Theorem 8.2, Corollary 8.3, and Example 8.4.}
Consider the fixed scalar equation $x^2+1=0$ in $\ZZ/p^m\ZZ$. The following are exact counts for five primes and three depths. The coefficient rings and the scalar nature of the unknown are fixed throughout this page.
\vfig{08_prime_forests}
\head{Where failure occurs matters}
At $p=3$ and $p=7$, no residual root exists. At $p=2$, the residual root exists but dies immediately: every odd square is one modulo four. At $p=5$ and $p=13$, the two residual roots are simple, and the usual one-digit equation gives one child above each root at every depth.

At five, one branch begins $2,7,57,182$; the other is its negative modulo each power of five. Arbitrarily many digits in either branch cannot supply the absent scalar root over $\ZZ_3$. Nor can one root modulo two satisfy the hypothesis that every finite level is nonempty.

The paper's system-completion theorem uses a supplied common algebraic spine, finite nonempty levels at every coefficient place, coherent reductions, and closed local conditions. The scalar example isolates why these are substantive requirements. It is not a disproof of a companion theorem whose coefficients may be enlarged or whose objects are higher-dimensional representations.
\takeaway{A residual root, roots at all depths, and roots at every required coefficient place are three distinct tests.}

\vpage{9}{Why finiteness enters compactness}{AN INFINITE-LEVEL COUNTERMODEL}{Main text: Theorem 3.3 and its finite-branching hypotheses.}
Let $X_m=\{n\in\ZZ:n\ge m\}$, with the inclusion $X_{m+1}\hookrightarrow X_m$ as transition map. Every level is nonempty. Every vertex has at most one child. Yet there is no compatible infinite branch.
\vfig{09_infinite_tail}
\head{The missing hypothesis is not surjectivity}
Compatibility would force all coordinates of a branch to be one fixed integer $n$. Membership at every level would then require $n\ge m$ for every $m$, which is impossible. Each vertical column in the drawing eventually ends. A finite window always contains columns extending beyond it, but no single column survives all windows.

There is no contradiction with finite-branch compactness: these levels are infinite. Attaching a common artificial root would give that root infinitely many children. The finitely branching rooted-tree argument therefore does not apply.

In the paper, finite nonempty levels and actual total reduction maps provide compactness even when some vertices die and some transition maps are not onto. One can retain only vertices with descendants at arbitrarily large depths. Finiteness makes that surviving part nonempty and permits compatible choices.
\[
X_m\ne\varnothing\ \text{for all }m
\quad\not\Longrightarrow\quad
\varprojlim_m X_m\ne\varnothing
\quad\text{without the compactness hypotheses.}
\]
\takeaway{Arbitrarily long finite paths are not automatically one infinite path. The finite-level hypothesis prevents escape to new roots.}

\vpage{10}{See cycles, boundaries, and classes}{A TWO-PACKET COCHAIN COMPLEX}{Main text: Theorem 4.2 and the finite-complex example.}
Over $\FF_5$, use bases $(e_0,e_1)$, $(u_0,u_1,u_2,u_3)$, and $(f_0,f_1)$ in degrees zero, one, and two. The only nonzero differential values are $d_0e_0=u_0$ and $d_1u_1=f_0$. Consequently $d_1d_0=0$.
\vfig{10_cohomology}
\head{The quotient is the important step}
In degree one, the cycles are $\Span(u_0,u_2,u_3)$, but $u_0$ is a boundary. Thus the cohomology is represented by $u_2,u_3$, not by every vector killed by $d_1$. In the other degrees the surviving classes are $e_1$ and $f_1$.
\[
\dim(H^0,H^1,H^2)=(1,2,1).
\]
Packet A contains $e_0$, $u_0,u_1$, and $f_0$. It is exact and contributes $(0,0,0)$. Packet B contains $e_1$, $u_2,u_3$, and $f_1$, with zero differentials, and contributes $(1,2,1)$. Adding the packet cohomologies reconstructs the total.

The projectors must be specified in every degree and commute with the differential. Their existence is not a consequence of the printed dimensions. Identifying this finite complex with the actual continuous cohomology of an arithmetic deformation problem would also require a separate comparison.
\takeaway{Resolve the differential first, then pass from cycles to classes. Counting packet dimensions alone does not compute cohomology.}

\vpage{11}{An arbitrary cut need not descend}{THE CHAIN-MAP CONDITION}{Main text: Theorem 4.2 and the scope of the cohomology splitting.}
Retain the previous complex. The valid A-packet projectors are $P^0=\diag(1,0)$, $P^1=\diag(1,1,0,0)$, and $P^2=\diag(1,0)$. They satisfy $d_rP^r=P^{r+1}d_r$.
\vfig{11_chain_defect}
\head{A boundary can cross an incompatible cut}
Replace the degree-one projector by $Q^1=\diag(0,1,0,0)$ while keeping $P^0$. Then
\[
d_0P^0e_0=u_0,
\qquad Q^1d_0e_0=0,
\qquad(d_0P^0-Q^1d_0)e_0=u_0\ne0.
\]
The selected degree-zero component is sent outside the selected degree-one component. There is no subcomplex with these cuts, so the cohomology packet theorem cannot be applied to them. Both matrices remain perfectly valid idempotents; the failure is compatibility with the differential.

This example separates two kinds of obstruction. The nonzero matrix above obstructs this proposed decomposition. A nonzero class in $H^2$ may obstruct a lift only after the deformation problem and its obstruction map have been identified. One should not equate every failed projector commutator with an arithmetic obstruction class.

The exact checks verify both commuting diagrams for the valid family and reject the altered family. They also enumerate the finite kernels and images, so the dimension comparison is not inferred from the drawing.
\takeaway{A projector on each vector space is not yet a projector on the complex. The differential must respect the chosen packets.}

\vpage{12}{How many starting seeds are necessary?}{DIRECTED TRANSPORT}{Main text: Theorem 7.3 and the seven-root example.}
The arrows below are the entire supplied transport graph. Vertices joined by paths in both directions form strongly connected components. The components $\{A,B\}$, $\{D,E\}$, and $\{G\}$ have no incoming edge from outside themselves.
\vfig{12_source_scc}
\head{A lower bound and a construction}
Any seed set reaching every vertex must meet all three source components. Otherwise there is no path into the omitted source. Conversely, one representative from each source reaches $C$ and $F$ as well. The directed seed rank is therefore three.

There are precisely four minimum seed sets:
\[
\{A,D,G\},\quad\{A,E,G\},\quad
\{B,D,G\},\quad\{B,E,G\}.
\]
The example checks this by exhaustive enumeration of all $2^7$ subsets, independently of the strongly connected component calculation. Adding an admissible edge cannot increase the required seed count. In particular, adding $F\to G$ makes $G$ reachable from the other sources, reducing the count to two.

The graph theorem is exact for the graph supplied. It does not establish that a proposed twist, induction, or arithmetic transfer preserves a determinant, local condition, or occupancy theorem. Those facts justify the arrows before graph optimization begins.
\takeaway{Count source components, not all components and not all vertices. But do not add an arithmetic arrow merely because it would improve the count.}

\vpage{13}{A transport must work at every depth}{ONE REVERSIBLE EXAMPLE}{Main text: Definition 7.1 and the reduction-compatibility requirement.}
For the two roots of $x^2+1=0$ modulo powers of five, negation is an actual reversible transport. It preserves the equation, is its own inverse, and commutes with reduction.
\vfig{13_transport}
\head{One successful input is not a natural transformation}
At every depth $m$,
\[
N_m(x)=-x\pmod{5^m},\qquad
r_{m+1,m}N_{m+1}=N_mr_{m+1,m}.
\]
The two compatible branches beginning at $2$ and $3$ are exchanged. Because $2\not\equiv3\pmod5$, this transport changes the residual root. It is horizontal rather than a positive-depth correction.

Translation by one is a useful false control. Modulo five it takes the root $2$ to the root $3$, but takes the root $3$ to $4$, which is not a root. A single successful arrow between selected points does not define an operation on the whole deformation problem.

Even negation is admissible only for the equation and conditions stated here. An extra condition that fixes the residual root at $2$ would not be preserved. Likewise, a representation-theoretic dual or twist needs its own compatibility checks for the chosen local data and carrier.
\takeaway{A transport is a family of maps respecting equations and reductions. A convenient pair of endpoints is not enough.}
\captiontext{Both branches are computed by exact one-digit lifting. Negation is checked on all roots at the displayed depths; its universal reduction identity follows directly from ring homomorphism laws.}

\vpage{14}{Common coefficients, different residues}{THE ARITHMETIC SPINE}{Main text: Definition 8.1 and Theorem 8.2.}
A common spine records algebraic data before reducing at a coefficient prime. For a small model, take the one polynomial
\[
P(T)=T^2-5T+6=(T-2)(T-3)\in\ZZ[T].
\]
Its reductions at five and seven look different, but they come from the same coefficients.
\vfig{14_spine}
\head{Agreement means descent from supplied data}
Modulo five the polynomial is $T^2+1$; modulo seven it is $T^2+2T+6$. The roots $2$ and $3$ are the reductions of the same integer roots in both places. The coefficient lists are not required to be equal as raw residues in different rings.

The paper's common spine is larger: it includes characteristic polynomials at the permitted arithmetic places and the specified determinant, weight, and local labels. Requiring every finite member to match those reductions is a compatibility condition on existing objects. It is not a construction of the objects.

Finite agreement also leaves ambiguity. Adding $35$ to an integer coefficient changes neither its residue modulo five nor its residue modulo seven. Thus two finite coefficient observations do not determine a unique global polynomial without a height bound or further exact information. Conversely, one arbitrarily chosen family of local objects need not descend to an algebraic spine.
\takeaway{The common spine is input shared across coefficient places. It is not created by comparing unrelated residue labels after the fact.}

\vpage{15}{A finite trace matrix can separate types}{CHARACTERS BEFORE ARITHMETIC PRIMES}{Main text: Theorem 9.1; latest continuation: arithmetic-system separation.}
Use the three irreducible characters of $S_3$: trivial, sign, and standard. On the identity, a transposition, and a three-cycle, their values form the matrix below. Rows are samples; columns are constituent types.
\vfig{15_separator}
\head{What full rank buys}
Since $\det C=6$, these samples distinguish every rational multiplicity vector. A target with two trivial constituents, no sign constituent, and one standard constituent has
\[
\bm m=\mat{2\\0\\1},\qquad
\bm b=C\bm m=\mat{4\\2\\1}.
\]
Solving $C\bm m=\bm b$ recovers the unique rational vector. The standard character is not introduced as an arbitrary list: the checks construct matrices for a transposition and a three-cycle, verify their group relations, and obtain their traces.

This finite example explains the linear algebra after separating samples are available. For the arithmetic theorem, the samples must be genuine unramified Frobenius classes with the required compatibility and allowed-prime conditions. Selecting abstract group elements does not construct such primes.

The latest continuation proves the final separator assembly from its explicit arithmetic-coverage proposition. The general prescribed-Frobenius occurrence proposition remains unproved in that release.
\takeaway{Invertibility answers reconstruction from these samples. Arithmetic occurrence answers whether the required samples can be taken at actual primes.}

\vpage{16}{A solution is not yet a summand}{INTEGRAL, NONNEGATIVE, WITHIN CAPACITY}{Main text: Theorem 9.2 and its multiplicity-recovery example.}
Keep the same separating matrix $C$ and suppose the available carrier has capacities $\bm c=(2,1,1)$. A recovered vector describes a genuine sub-sum of that carrier only if each entry is an integer and $0\le m_i\le c_i$.
\vfig{16_multiplicities}
\head{Three independent rejection tests}
The trace vector $(1,0,0)$ gives $(1/6,1/6,1/3)$, so it fails integrality. The vector $(2,-2,-1)$ gives $(-1,1,1)$, which is integral but not nonnegative. The vector $(5,3,2)$ gives $(3,0,1)$, which exceeds the available capacities. In each case the linear equation itself has a unique rational solution.

The accepted example $\bm b=(4,2,1)$ recovers $(2,0,1)$. When compatible constituent families are already supplied, the same admissible multiplicities define the corresponding direct sum at every coefficient place. This is the finite algebraic core of summand transport.

It does not supply a compatible family for an arbitrary object, construct all local attachments, or ensure that a proposed carrier contains the desired representation. The latest Paper X ledger leaves those remaining arithmetic comparisons visible in partial Claim 9.2.
\takeaway{Invert the trace matrix, then check the representation-theoretic constraints. A unique vector in a field need not be a realizable multiplicity vector.}
\captiontext{All accepted and rejected vectors are solved over exact rationals. The code rejects the wrong cases by their stated reason, not by a floating-point tolerance.}

\vpage{17}{Separation can consume a digit}{EXACT PRECISION LOSS}{Main text: The finite trace-recovery discussion in §9.}
The character matrix has determinant six, so its inverse contains a factor of two in the denominator. That arithmetic fact is visible as a loss of one binary digit in trace reconstruction.
\vfig{17_precision}
\head{A concrete ambiguity modulo eight}
The multiplicity difference $\delta\bm m=(4,4,0)$ is nonzero modulo eight, but
\[
C\,\delta\bm m=(8,0,8)\equiv0\pmod8.
\]
Thus traces known modulo eight do not distinguish those two residue vectors. The same trace difference is nonzero modulo sixteen, so one more trace digit resolves this ambiguity.

More generally, the kernel of $C$ modulo $2^n$ has exactly two elements:
\[
(0,0,0),\qquad(2^{n-1},2^{n-1},0).
\]
To see this, the second equation gives $m_1=m_2$; the third gives $m_3=2m_1$; the first then gives $6m_1=0$. Three is a unit modulo $2^n$, leaving precisely the displayed possibilities.

Conversely, $C\delta\bm m\equiv0\pmod{2^{n+1}}$ implies $6\delta\bm m\equiv0\pmod{2^{n+1}}$ by the adjugate identity, hence $\delta\bm m\equiv0\pmod{2^n}$. Capacity bounds can remove ambiguity earlier in a particular carrier, but do not change this universal linear precision calculation.
\takeaway{Full rank over characteristic zero is not the same as invertibility over every residue ring. Denominator valuations determine the needed extra precision.}

\vpage{18}{Pairwise occupancy can still fail jointly}{EXACT JOINT PROJECTORS}{Main text: Theorem 10.1 and the finite occupancy models.}
Let a six-dimensional rational carrier have the joint eigenvalue labels shown below for three commuting diagonal operators $T_1,T_2,T_3$. Each coordinate is zero or one. Selecting a zero uses $I-T_i$; selecting a one uses $T_i$.
\vfig{18_occupancy}
\head{The full conjunction is a different test}
For target $(0,0,0)$, successive conditions reduce the ranks from $6$ to $3$, then $2$, then $1$. The joint projector is nonzero, so that label is occupied in this carrier.

For target $(0,1,1)$, every pair of conditions is occupied: $(0,1,0)$ witnesses the first two, $(0,0,1)$ the first and third, and $(1,1,1)$ the second and third. But the triple itself is absent. Multiplying all three selectors gives the zero matrix.
\[
E_{\bm a}=\prod_{i=1}^{3}
\begin{cases}I-T_i,&a_i=0,\\T_i,&a_i=1.\end{cases}
\qquad \Tr(E_{\bm a})=\rank(E_{\bm a}).
\]
The checks compute all eight joint projectors, their pairwise orthogonality, and their sum. The sum is the identity even though two labels are absent: those two projectors are simply zero.
\takeaway{A formula for a projector does not force its image to be nonzero. Occupancy must be tested on the complete set of required conditions.}

\vpage{19}{The hypotheses do real work}{TRACE AND SEMISIMPLICITY}{Main text: Theorem 10.1 and the checked false controls following it.}
The finite occupancy theorem uses a split semisimple commuting family in characteristic zero. Two tiny matrices show why neither the coefficient condition nor semisimplicity should disappear from the statement.
\vfig{19_hypotheses}
\head{Trace zero need not mean an empty image}
On $\FF_5^5$, the identity is an idempotent of rank five, but its field-valued trace is $5=0$ in $\FF_5$. Rank and localization still detect its nonzero image. A characteristic-zero integer trace cannot be silently transferred to this setting.
\head{An eigenvalue label need not select ordinary eigenvectors}
For the nilpotent matrix $N=\left(\begin{smallmatrix}0&1\\0&0\end{smallmatrix}\right)$, the only eigenvalue is zero. A Lagrange selector that multiplies over all other eigenvalues is an empty product, hence $I$. Its image has dimension two, while $\ker N$ has dimension one. The selector does not project onto the ordinary eigenspace.

The whole two-dimensional space is the generalized zero-eigenvalue packet. Paper XII's Artin-packet treatment is designed to retain such nilpotent structure. That is a different target from Paper X's split semisimple joint-eigenspace formula.
\takeaway{A correct theorem in one coefficient category is not a theorem in all categories. Changing from ordinary eigenspaces to generalized packets also changes the object being detected.}

\vpage{20}{A condition must survive the limit}{CLOSEDNESS IS NOT DECORATION}{Main text: Theorems 8.2 and 11.1; the inverse-limit closure hypotheses.}
In $\ZZ_3$, every term of the sequence $x_k=3^k$ is nonzero. Yet $x_k\to0$ in the three-adic topology. The property ``nonzero'' is not closed, so a limit of objects satisfying it can lose that property.
\vfig{20_closedness}
\head{Polynomial conditions behave differently}
For a continuous polynomial $f$, if a convergent sequence satisfies $f(x_k)=0$ for every $k$, then continuity gives $f(x)=0$ at its limit. More generally, a compatible residue branch satisfying $f(x_m)\equiv0\pmod{3^m}$ reconstructs an exact root once the complete ring and its reduction maps have been identified.

The sequence $3^k$ here is \emph{not} presented as a compatible branch of nonzero residues modulo $3^k$: those particular residues are zero. It is a convergent sequence of nonzero elements of the complete ring. The distinction prevents an incorrect finite-level interpretation of the picture.

In an arithmetic application, local conditions may be given by quotient equations, integral structures, or other data. The completion theorem requires a proof that the chosen condition is preserved by the limit under construction. The word ``admissible'' does not itself provide that proof.
\takeaway{Compatibility reconstructs an object. Closedness is what carries the declared equations or conditions to that object.}
\captiontext{The finite diagram displays the retreat of the first nonzero digit. The convergence assertion follows from $|3^k|_3=3^{-k}$, not from extrapolating a finite numerical plot.}

\vpage{21}{The sign of a norm is arithmetic data}{ELEMENTS VERSUS IDEALS}{Main text: §9, arithmetic residue-counting continuation; signed-norm comparison.}
Even in the number field $\QQ$, congruent generators need not have congruent absolute norms. The integers $2$ and $-3$ are congruent modulo five. Their signed field norms are congruent, but their absolute values are not.
\vfig{21_signed_norm}
\head{The smallest counterexample is already decisive}
For $K=\QQ$,
\[
N_{K/\QQ}(2)=2,\qquad N_{K/\QQ}(-3)=-3\equiv2\pmod5,
\]
whereas
\[
N((2))=2,\qquad N((-3))=3\not\equiv2\pmod5.
\]
Here $N((a))$ is the positive index of the ideal $(a)$, equal to $|a|$. The determinant norm of an element is signed. Taking absolute values is not a ring operation and does not commute with reduction in this example.

The latest continuation counts the two real norm signs separately before reassembling its generator-to-ideal sums. Its exact comparison also includes the roots-of-unity multiplicity and the norm of the chosen inverse-class representative. Omitting those factors would change the arithmetic sum, not merely its normalization.

The finite-field norm on the following pages is a determinant in a finite algebra. It has no ordering or absolute-value sign. Its role is to control residue characters before the separate real-lattice counting argument.
\takeaway{Do not infer an ideal-norm residue from a generator congruence without accounting for the sign in the element-to-ideal comparison.}

\vpage{22}{A norm is a determinant you can compute}{THE QUOTIENT ALGEBRA $\FF_{25}$}{Main text: §9; latest continuation: actual quotient determinant norms.}
Use $A=\FF_5[u]/(u^2-2)$. Since two is not a square modulo five, this quotient is a field with 25 elements. Every element is uniquely $a+bu$, with $a,b\in\FF_5$.
\vfig{22_norm_table}
\head{Coordinates, multiplication, and norm agree}
In the basis $(1,u)$, multiplication by $a+bu$ has matrix
\[
M_{a,b}=\mat{a&2b\\b&a},
\qquad N(a+bu)=\det M_{a,b}=a^2-2b^2\pmod5.
\]
The table shows every norm value. Only the zero element has norm zero. Each nonzero norm value has exactly six preimages. This makes the later character cancellation visible without estimating any sums.

The checks enumerate all 25 matrices, verify the determinant formula, and verify norm multiplicativity on every pair of field elements. They also change to a different invertible basis and check that conjugating the multiplication matrix leaves its determinant unchanged.

This model is the reduction of the quadratic integer algebra $\ZZ[\sqrt2]$ modulo five. Its finite arithmetic illustrates the determinant construction used in the continuation; it is not a computation of all ideal classes or a general ideal-counting theorem.
\takeaway{A quotient norm is built from the actual multiplication operator. Its basis independence is a determinant identity, not a choice of convenient coordinates.}

\vpage{23}{Permutation forces a zero mean}{A NONSCALAR UNIT WITNESS}{Main text: §9; latest continuation: norm-residue lattice mean cancellation.}
Let $\chi$ be the quadratic character modulo five, extended by $\chi(0)=0$. In the field on the preceding page, choose $v=2u$. Its norm is $N(v)=2$, so $\chi(N(v))=-1$.
\vfig{23_mean_cancellation}
\head{The whole sum is paired with its negative}
Multiplication by $v$ permutes the 25 elements. Norm multiplicativity gives
\[
\sum_{z\in A}\chi(N(z))
=\sum_{z\in A}\chi(N(vz))
=-\sum_{z\in A}\chi(N(z)).
\]
The sum is an integer, so it is zero. The origin contributes zero. The other 24 elements fall into three cycles of length eight under multiplication by $v$; the character signs alternate around each cycle.

The same reasoning works with a complex root of unity $\chi(N(v))\ne1$: subtracting the two equal sums gives $(1-\chi(N(v)))S=0$. The essential inputs are an actual permutation, a multiplicative norm, and a nontrivial character value. They are not consequences of drawing an alternating cycle.

In the arithmetic continuation the quotient is the actual additive module $J/mJ$ associated with an ideal representative. This finite-field example isolates the cancellation mechanism while keeping that broader module and lattice comparison separate.
\takeaway{The witness is one unit that changes the character value. Once it exists, the permutation cancels the entire finite mean.}

\vpage{24}{A nontrivial character may still be exceptional}{THE EXACT SCALAR TEST}{Main text: §9; latest status: scalar-power and nonscalar-witness boundaries.}
A scalar $a\in\FF_5^\times$ acts on a degree-$d$ algebra with determinant $a^d$. Therefore the scalar cancellation route tests $\chi(a^d)$, not merely $\chi(a)$.
\vfig{24_scalar_exception}
\head{A failed sufficient test is not an impossibility theorem}
Let $\chi_4(2)=i$ generate the character group of $\FF_5^\times$, and write $\chi_k=\chi_4^k$. Since $2$ generates the unit group,
\[
\exists a:\chi_k(a^d)\ne1
\quad\Longleftrightarrow\quad kd\not\equiv0\pmod4.
\]
For the quadratic character $k=2$ and degree $d=2$, every scalar test is trivial. Yet page \pagerefV{23} supplies the nonscalar witness $2u$, proving cancellation in that specific algebra. Scalar exceptional does not mean exceptional for every possible quotient unit.

Conversely, the quartic character $k=1$ works in degree two even though $\gcd(2,4)\ne1$. Thus the simpler degree--totient coprimality criterion is sufficient, not necessary. With degree three, every nontrivial character passes because exponentiation by three permutes the units.

The latest continuation states its public hypothesis in the weaker exact scalar-power form. It does not claim to construct a nonscalar witness for every scalar-exceptional arithmetic character.
\takeaway{Test the character after the degree-power map. Nontriviality before that map can disappear after it.}

\vpage{25}{Avoid finitely many primes constructively}{A CRT MODEL FOR REPRESENTATIVE CHOICE}{Main text: §9; latest continuation: finite-prime avoidance and automatic representatives.}
A representative needed by an arithmetic argument often has to avoid a prescribed finite collection of primes. A small integer version asks for
\[
x\equiv1\pmod2,\qquad x\equiv2\pmod3,\qquad x\equiv4\pmod5.
\]
There is exactly one solution modulo 30: $x=29$.
\vfig{25_avoidance}
\head{Avoidance and ideal class are different constraints}
The selected residues are nonzero modulo each forbidden prime, so $\gcd(29,30)=1$. The eight units modulo 30 are displayed separately. They all avoid the primes, but only one also satisfies the three specific residue prescriptions.

The continuation proves an ideal-theoretic analogue with more structure: in every prescribed inverse ideal class it constructs a nonzero integral representative whose norm is coprime to the nonzero modulus. That removes a representative-choice assumption from the public norm-character bound. It does not remove the remaining character, lattice, or Frobenius-occurrence obligations.

The integer CRT calculation is an illustration, not a proof of the ideal-class representative theorem. In particular, enumerating integers cannot impose a nontrivial ideal class. The full construction uses prime-ideal avoidance and its source comparison to the class group.
\takeaway{A finite avoidance argument supplies an admissible representative. It does not by itself make the arithmetic object being counted, or the desired Frobenius class, occur.}
\captiontext{All 30 residue classes are tested exactly. The unique CRT result and the unit count are independent finite outputs of the check script.}

\vpage{26}{Cancellation leaves a boundary problem}{A FULL LATTICE-BOX CALCULATION}{Main text: §9; latest continuation: additive residue counts and norm-character bounds.}
Extend the function $f(a,b)=\chi(a^2-2b^2\bmod5)$ periodically to $\ZZ^2$. Every complete $5\times5$ residue tile has sum zero by the finite-field calculation. Only incomplete tiles can contribute to a box sum.
\vfig{26_boundary_count}
\head{Separate bulk cancellation from boundary size}
For $0\le a,b<12$, four complete tiles fill a $10\times10$ square and contribute zero. The remaining 44 points lie in the two boundary strips. Their exact sum is one. Since $|f|\le1$, a bound of 44 follows without needing that final cancellation.

For a general positive integer $L$, write $L=5q+r$, $0\le r<5$. The complete tiled square has side $L-r$, so
\[
|S_L|\le L^2-(L-r)^2=2Lr-r^2\le8L.
\]
The area is $X=L^2$, making this a square-root-in-area bound for this square-box model. The proof is the exact tile decomposition; checking many $L$ values is only a regression check of the implementation.

Number-field ideal sums involve different regions, lattice indices, signs, and generator multiplicities. Their $X^{1-1/d}$ exponent cannot be justified simply by relabeling this square picture as an ideal-norm ball.
\takeaway{A zero finite mean removes the bulk. The remaining estimate depends on the actual geometry of the region and its boundary.}

\vpage{27}{What a power-saving bound buys analytically}{A MELLIN ENDPOINT MODEL}{Main text: §9; latest continuation: summatory measurability and first-order Mellin control.}
Use the explicit envelope $A(x)=x^{1/2}$ for $x\ge1$. It is not an ideal-counting function. It is a model showing why a power saving can make a Mellin expression regular at $s=1$.
\vfig{27_mellin}
\head{Compute the endpoint rather than just naming it}
For real $s>1/2$,
\[
I(s)=s\int_1^\infty A(x)x^{-s-1}\,\dd x
=\frac{s}{s-1/2}.
\]
At $s=1+h$ this gives
\[
I(1+h)-I(1)=-\frac{h}{h+1/2},
\qquad |I(1+h)-I(1)|\le4|h|
\quad (|h|\le1/4).
\]
The derivative at one is $-2$. The logarithmic moment is finite as well: $\int_1^\infty x^{-3/2}\log x\,\dd x=4$. These are explicit calculations; the drawn curve does not establish differentiability or convergence.

The continuation performs the corresponding analytic steps for its actual summatory functions, with their explicit measurability and growth hypotheses. The arithmetic Euler-product identity, nonzero pole residues, and the required prescribed-prime divergence remain separate. Regularity of one Mellin transform does not construct all of those inputs.
\takeaway{An analytic estimate controls a function already tied to arithmetic data. It is not a substitute for proving that tie.}

\vpage{28}{Finite splitting data are not a density theorem}{THE REMAINING FROBENIUS STEP}{Main text: Theorem 9.1. Latest continuation: GeneralArithmeticCoverage and ArithmeticSystemSeparation.}
For the quadratic polynomial $X^2+1$, reduction modulo an odd prime either gives two roots or no roots. This finite scan displays both behaviors. It is a useful model for the distinction between splitting and nonsplitting classes.
\vfig{28_frobenius_boundary}
\head{Existence must survive arbitrary exclusions}
The continuation's missing general coverage proposition asks for an actual unramified rational prime in a prescribed conjugacy class, outside any prescribed finite forbidden set, together with the corresponding absolute Frobenius witness. A list of small primes does not establish that quantifier.

The finite character argument on page \pagerefV{15} also needs samples that jointly separate the constituents. Repeating the identity-class row $(1,1,2)$ of the $S_3$ table gives rank one no matter how often it is repeated. Many samples are not necessarily informative samples.

The latest checked separator theorem composes actual arithmetic coverage with the finite linear-algebra construction. Thus the final assembly is available once the coverage proposition is supplied. The general occurrence theorem itself is not a field in that finite matrix example and is not certified by the displayed scan.
\takeaway{A separating group element, a witnessed Frobenius at one prime, and a new prime after every finite exclusion are distinct levels of evidence.}
\captiontext{The scan exhaustively checks scalar roots for the shown primes. No density, infinitude, or general prescribed-conjugacy statement is inferred from it.}

\vpage{29}{Changing coefficients changes the problem}{SCALAR ABSENCE VERSUS REPRESENTATION EXISTENCE}{Main text: Example 8.4 and the declared coefficient categories in §§8--9.}
Over $\FF_3$, no scalar solves $x^2+1=0$. But the matrix
\[
A=\mat{0&-1\\1&0}\in M_2(\FF_3)
\quad\text{satisfies}\quad A^2+I=0.
\]
And over $\FF_9=\FF_3[u]/(u^2+1)$, the scalars $u$ and $-u$ solve the same polynomial.
\vfig{29_coefficients}
\head{Three outputs with different types}
The failed scalar problem over $\FF_3$, the successful matrix problem over $\FF_3$, and the successful scalar problem over $\FF_9$ are not contradictory answers. They have different allowed objects. The matrix is multiplication by $u$ in the $\FF_3$-basis $(1,u)$ of the extension field.

This also has a characteristic-zero analogue. The matrix with the same entries gives a rational two-dimensional representation of a cyclic group of order four. Its scalar eigenvalues require a field containing a square root of $-1$. No one-dimensional rational eigencharacter is created by writing the two-dimensional matrix.

Accordingly, the scalar companion failure on page \pagerefV{8} is a counterexample only to the fixed scalar problem. It is not a counterexample to a representation problem that permits coefficient extension or a different dimension. Those permissions must be declared before comparing fibers across primes.
\takeaway{An empty fiber is meaningful only after its coefficient ring and object type are fixed. Do not repair an empty fiber by silently changing either.}

\vpage{30}{Read the evidence at its actual scope}{FINAL READER'S CHECKLIST}{Main text: cumulative formalization notes and Claim 11.1. Latest continuation: CURRENT\_STATUS.md.}
The selected source history has three distinct layers: the amended Block 27 manuscript, the subsequent residue-counting release, and the later scalar/Mellin plus arithmetic-separator continuation. The most recent status file governs the formal coverage, not the older manuscript's prose ledger.
\par\noindent\begin{tabularx}{\linewidth}{@{}p{.19\linewidth}Y@{}}\toprule
Status & What is retained\\\midrule
\textbf{17 checked} & The source ledger's completed claims, under their recorded hypotheses. Conditional completion results retain actual finite levels, reductions, local relations, and inverse-limit closure premises.\\
\textbf{9.1 partial} & Automatic representatives, norm-character estimates, Mellin control, and the final separator assembly are checked in saved records. General prescribed Frobenius occurrence remains unproved; the analytic alternative retains its arithmetic product, residue, and divergence inputs.\\
\textbf{9.2 partial} & Supplied-data summands are constructed. Actual Hodge--Tate/Weil--Deligne attachments, coefficient comparisons, and remaining general recovery are not all finished.\\
\textbf{11.1 partial} & The conditional core remains checked. Required horizontal transports and the remaining local preservation and recovery obligations are unfinished.\\\bottomrule
\end{tabularx}
\head{What was added here}
The appendix supplies reproducible finite examples and their exact Python checks. Geometry may be drawn using decimal positions, but no plotted pixel decides a congruence, rank, determinant, root, or character sum. The original Lean sources and recorded audits are preserved. This editorial addition does not run Lean, introduce new axioms, or promote a partial claim.
\head{The questions to carry into an application}
Specify the residual object and allowed coefficients. Prove that the finite levels and their maps model the intended problem. Check closedness before passing to a limit. Justify each transport. Establish the arithmetic samples before recovering multiplicities. Finally, exhibit nonzero occupancy in the required carrier rather than infer it from a projector formula.
\takeaway{The paper's strength is separation: it identifies which construction answers each question, and prevents one completed layer from being used as evidence for a different missing layer.}
\captiontext{Source provenance, replay commands, dependency pins, exact checks, and unchanged-input hashes accompany this appendix. No historical-priority claim or unrestricted companion/automorphic existence theorem is added.}
